\documentclass[11pt]{article}
\usepackage[margin=1in]{geometry}
\usepackage{amsmath,amssymb,amsthm}
\usepackage{mathtools}
\usepackage{hyperref}
\usepackage{booktabs}
\usepackage{array}
\usepackage{algorithm}
\usepackage{algpseudocode}
\usepackage{xcolor}

\algrenewcommand\algorithmicrequire{\textbf{Input:}}
\algrenewcommand\algorithmicensure{\textbf{Output:}}
\algrenewcommand\algorithmiccomment[1]{\hfill\(\triangleright\) {\footnotesize #1}}

\newtheorem{theorem}{Theorem}
\newtheorem{lemma}{Lemma}
\newtheorem{definition}{Definition}
\newtheorem{remark}{Remark}

\newcommand{\Gone}{\mathbb{G}_1}
\newcommand{\Gtwo}{\mathbb{G}_2}
\newcommand{\Gt}{\mathbb{G}_T}
\newcommand{\Fr}{\mathbb{F}_r}
\newcommand{\msg}{\mathrm{msg}}
\newcommand{\Enc}{\mathsf{Enc}}
\newcommand{\Dec}{\mathsf{Dec}}
\newcommand{\Adv}{\mathrm{Adv}}
\newcommand{\negl}{\mathrm{negl}}

\title{Partial-Binding Witness Encryption over Groth16}
\author{Stephen Duan\\\texttt{stephen.d@zkm.io}}
\date{2026-04-19}

\begin{document}
\maketitle

\begin{abstract}
BABE's witness encryption for Groth16 binds the ciphertext to the
full public-input vector at encryption time, precluding protocols
where some inputs are determined after encryption. We introduce
\emph{partial-binding WE}, which binds only a designated static
prefix $x_S$ and lets the remaining indices $x_D$ be committed at
dispute time via a \emph{Dual-Scalar Garbled Circuit} (DSGC) that
outputs both $r \cdot \pi_1$ and $r \cdot P_D = r \sum_{j \in D} x_j L_j$.
The DSGC structure binds $x_D$ proactively to the specific proof
used in decryption, closing a substitution attack that any
non-binding construction admits. Overhead is $|D|$ additional
fixed-base scalar multiplications inside the garbled circuit.
Security reduces to basic BABE plus GC-security and Lamport
one-wayness. The scheme retains BABE's property of requiring no
on-chain Groth16 verifier.
\end{abstract}

% ============================================================
\section{Introduction}

Witness encryption (WE) over Groth16~\cite{groth16}, as instantiated
by BABE~\cite{babe2026}, produces a ciphertext that decrypts to a
hidden message iff the decryptor supplies a valid Groth16 proof
$\pi$ for a fixed statement $(vk, x)$. The encryption procedure
$\Enc(vk, x, \msg)$ consumes the entire public-input vector $x$ at
setup time.

This rigidity is a deployment obstacle. Consider a protocol in which
a Verifier encrypts at time $T_0$ and a Prover decrypts at time
$T_1 > T_0$, where some of the public-input values are determined by
events occurring in $[T_0, T_1]$ (e.g.\ blockchain state evolving
between setup and proof time). At $T_0$ the Verifier has access only
to a prefix $x_S$ of $x$; the suffix $x_D$ is revealed at $T_1$.
Basic BABE cannot encrypt under these constraints.

\subsection{Contribution}

We present partial-binding WE, a modification of BABE in which
$\Enc$ takes only $(vk, x_S, S, D, \msg)$ as input and produces an
augmented ciphertext. Decryption succeeds for any valid
$(x_S, x_D, \pi)$ with $x_S$ matching the committed prefix. The
modification is:

\begin{itemize}
  \item \textbf{Encryption:} replace $Y = e(\alpha, \beta) \cdot e(vk_x, \gamma)$ with
    $Y_S = e(\alpha, \beta) \cdot e(P_S, \gamma)$ where
    $P_S = \sum_{j \in S} x_j L_j$; publish additional
    $ct_D = \{r \cdot L_j\}_{j \in D} \in \Gone^{|D|}$.
  \item \textbf{Decryption:} compute $r \cdot P_D$ from $x_D$ and
    $ct_D$, strip the dynamic factor $e(P_D, \gamma)^r$ from
    $r \cdot Y$ via bilinearity, unmask under $Y_S$.
\end{itemize}

Security reduces to BABE WE security plus DLog hardness in $\Gone$
(Lemma~\ref{lem:pbwe}). On-chain cost is preserved: decryption runs
off-chain in both schemes, and the only on-chain delta is $|D|$
additional Lamport/Winternitz commitments to pin $x_D$ via Bitcoin
script.

The construction exploits the linearity of Groth16's public-input
aggregation $vk_x = \sum_j x_j \cdot L_j$. SNARKs with non-linear
public-input handling do not directly admit this technique.

\subsection{Outline}

\S\ref{sec:prelim} establishes notation and recalls basic BABE.
\S\ref{sec:construction} presents our construction.
\S\ref{sec:security} gives the security analysis.
\S\ref{sec:cost} summarizes cost.

% ============================================================
\section{Preliminaries}\label{sec:prelim}

\subsection{Notation}

Let $e : \Gone \times \Gtwo \to \Gt$ denote a Type-III pairing on
BN254. The Groth16 verifying key is
$vk = (\alpha, \beta, \gamma, \delta, \{L_j\}_{j=0}^{n})$ where
$\alpha, \delta \in \Gone$ and $\beta, \gamma \in \Gtwo$ (we abuse
notation and treat $\delta$ as its $\Gtwo$ element throughout). A
proof is $\pi = (\pi_1, \pi_2, \pi_3) \in \Gone \times \Gtwo \times \Gone$.
The public-input vector is $x = (x_0, x_1, \ldots, x_n) \in \Fr^{n+1}$
with $x_0 = 1$. The verifier accepts iff
\begin{equation}
  e(\pi_1, \pi_2) \;=\; e(\alpha, \beta)\,\cdot\,e(vk_x, \gamma)\,\cdot\,e(\pi_3, \delta),
  \qquad vk_x \;:=\; \sum_{j=0}^{n} x_j \cdot L_j.
  \label{eq:g16}
\end{equation}

$H : \{0,1\}^* \to \{0,1\}^{256}$ is a collision-resistant hash. We
write $\Adv_{\mathrm{scheme}}(\lambda)$ for the advantage of a PPT
adversary at security parameter $\lambda$, and $\negl(\lambda)$ for
negligible functions.

\subsection{Witness encryption for Groth16}

\begin{definition}[WE for Groth16]
A witness encryption scheme for Groth16 relation
$\mathcal{R}_{vk} = \{(x, w) : \exists \pi = \mathsf{Prove}(vk, x, w)
\text{ accepts}\}$ is a pair of PPT algorithms $(\Enc, \Dec)$ such
that:
\begin{itemize}
  \item \emph{Correctness:} for every $(x, w) \in \mathcal{R}_{vk}$,
    $\Dec(\mathsf{Prove}(vk, x, w),\, \Enc(vk, x, \msg)) = \msg$.
  \item \emph{Soundness:} for every $x \notin \mathcal{L}_{vk}$ (i.e.,
    $\{w : (x, w) \in \mathcal{R}_{vk}\} = \emptyset$), and every
    $\msg \neq \msg'$,
    $\{\Enc(vk, x, \msg)\} \approx_c \{\Enc(vk, x, \msg')\}$.
\end{itemize}
\end{definition}

BABE~\cite{babe2026} instantiates $(\Enc, \Dec)$ with a pairing-based
construction; we recall it below.

\subsection{BABE ciphertext structure}

A BABE ciphertext per cut-and-choose instance is
$ct = (gc\_ct, adaptor, ct_2, ct_3)$ where
$gc\_ct$ is a garbled circuit computing $r \cdot \pi_1$, $adaptor$
binds its output labels to Bitcoin pre-signatures, $ct_2 = r \cdot \delta$,
and $ct_3 = \msg \oplus H(r \cdot Y)$ with
$Y = e(\alpha, \beta) \cdot e(vk_x, \gamma)$. The per-instance scalar
$r$ and mask $\msg$ are derived from a seed via a PRG. Cut-and-choose
over $N_{\mathrm{CC}}$ instances detects malformed garbling with
probability exponential in the opened subset size.

\begin{algorithm}[H]
\caption{BABE $\Enc(vk, x;\, I)$ \hfill{\footnotesize Verifier, off-chain, $T_0$}}
\label{alg:babe-enc}
\begin{algorithmic}[1]
\Require $vk$; full public inputs $x = (x_0, \ldots, x_n)$;
  instance $I = (r, \msg, \Delta, gc\_ct, adaptor)$
\Ensure $ct = (gc\_ct, adaptor, ct_2, ct_3)$
\State $ct_2 \gets r \cdot \delta$
\State $vk_x \gets \sum_{j=0}^{n} x_j \cdot L_j$ \Comment{requires every $x_j$}
\State $Y \gets e(\alpha, \beta) \cdot e(vk_x, \gamma)$
\State $ct_3 \gets \msg \oplus H(r \cdot Y)$
\State \Return $(gc\_ct, adaptor, ct_2, ct_3)$
\end{algorithmic}
\end{algorithm}

\begin{algorithm}[H]
\caption{BABE $\Dec(\pi, ct)$ \hfill{\footnotesize Prover, off-chain, $T_1$}}
\label{alg:babe-dec}
\begin{algorithmic}[1]
\Require proof $\pi = (\pi_1, \pi_2, \pi_3)$ for $(vk, x)$; $ct$
\Ensure $\msg$
\State $ct_1 \gets r \cdot \pi_1$ \Comment{via GC unlock; $r$ not revealed}
\State $r \cdot Y \gets e(ct_1, \pi_2) / e(\pi_3, ct_2)$ \Comment{by \eqref{eq:g16}}
\State $\msg \gets ct_3 \oplus H(r \cdot Y)$
\State \Return $\msg$
\end{algorithmic}
\end{algorithm}

\begin{remark}[On-chain surface]
Algorithms~\ref{alg:babe-enc} and \ref{alg:babe-dec} run off-chain.
Bitcoin script is used only for (i) hash commitments to
$(ct_2, ct_3)$ published during C\&C, (ii) gate-by-gate GC unlocking
at dispute time, (iii) a hashlock on $\msg$ gating the final UTXO
spend. No pairings or target-group arithmetic are evaluated on chain.
\end{remark}

% ============================================================
\section{Partial-binding witness encryption}\label{sec:construction}

\subsection{Problem statement}

\begin{definition}[Partial-binding WE]\label{def:pbwe}
Fix a partition $\{0, 1, \ldots, n\} = S \sqcup D$ at circuit
compile time with $0 \in S$. A partial-binding WE scheme for
$\mathcal{R}_{vk}$ with respect to $(S, D)$ is a pair
$(\Enc^\star, \Dec^\star)$ such that:
\begin{itemize}
  \item \emph{Prefix-only encryption:} $\Enc^\star$ takes input
    $(vk, x_S, S, D, \msg)$; it does not depend on $x_D$.
  \item \emph{Correctness:} for every $((x_S, x_D), w) \in \mathcal{R}_{vk}$
    with $\pi = \mathsf{Prove}(vk, (x_S, x_D), w)$, it holds that
    \[
      \Dec^\star\bigl(\pi,\,(x_S, x_D),\,\Enc^\star(vk, x_S, S, D, \msg)\bigr) = \msg.
    \]
  \item \emph{Prefix soundness:} for every $x_S$ such that no
    $(x_D, w)$ satisfies $\mathcal{R}_{vk}$, and every $\msg \neq \msg'$,
    \[
      \{\Enc^\star(vk, x_S, \cdot, \msg)\}
      \;\approx_c\;
      \{\Enc^\star(vk, x_S, \cdot, \msg')\}.
    \]
\end{itemize}
\end{definition}

Partial-binding does not constrain $x_D$: decryption succeeds for
any $(x_D, \pi)$ with $\pi$ valid for $(x_S, x_D)$. External
mechanisms (e.g.\ Lamport commitment plus Bitcoin-script equality)
must pin $x_D$ for enclosing protocols. We formalize this separation
in \S\ref{sec:security}.

\subsection{Dual-Scalar Garbled Circuit}\label{sec:dsgc}

Basic BABE's garbled circuit computes a single scalar multiplication
$r \cdot \pi_1$. We replace it with a \emph{Dual-Scalar Garbled Circuit}
(DSGC) that computes two scalar-mul outputs:
\[
  \mathsf{DSGC}(\pi_1,\ x_D) \;\longrightarrow\; (r \cdot \pi_1,\ r \cdot P_D),
\]
where $P_D = \sum_{j \in D} x_j \cdot L_j$ is the dynamic-prefix
contribution and the scaled bases $\{r \cdot L_j\}_{j \in D}$ are
baked into the circuit as garbling-time constants (the Verifier
knows both $r$ and the CRS $\{L_j\}$ at $T_0$).

The DSGC input-label layout reserves Lamport-committed wires for
both $\pi_1$ (as in basic BABE) and $x_D$ (new). At dispute, the
Prover Lamport-reveals both; BABE's soldering step routes those
reveals into the DSGC input wires, producing the two scalar-mul
outputs via the standard GC-evaluation path.

\paragraph{Why DSGC binds $x_D$ to the proof.}
The $r \cdot P_D$ output is uniquely determined by the
Lamport-committed $x_D$: it is not the Prover's arbitrary off-chain
choice. The bilinear cancellation in decryption
(Algorithm~\ref{alg:pb-dec}, line 4) then fails whenever the proof's
public input differs from the Lamport-committed $x_D$, because the
two $e(P_D, \gamma)^r$ factors no longer match. We formalize this in
\S\ref{sec:security}.

\paragraph{Lamport-$x_D$ soldering (required).}
The binding in Alg.~\ref{alg:dsgc} relies on two properties:
\begin{enumerate}
  \item \emph{GC privacy} hides the scaled bases
    $\{r \cdot L_j\}_{j \in D}$ baked into the DSGC as garbling-time
    constants. The evaluator learns neither individual constants nor
    intermediate wire values; only the final outputs
    $(r \cdot \pi_1, r \cdot P_D)$ for the specific inputs provided.
  \item \emph{Lamport-$x_D$ soldering} ties Lamport preimages of
    $x_D$ bits to the DSGC input-wire labels, so the only $x_D$ that
    can be fed into the DSGC is the on-chain-committed one. For each
    bit $(j, i)$ of $x_D$, the Verifier samples Lamport preimage pair
    $(s_{j,i}^0, s_{j,i}^1)$, publishes
    $(H(s_{j,i}^0), H(s_{j,i}^1))$, and derives DSGC input labels
    $L_{j,i}^b = \mathrm{KDF}(s_{j,i}^b)$ so that Lamport reveal
    uniquely determines the GC input label. This mirrors BABE's
    soldering of Lamport-$\pi_1$ to the GC's $\pi_1$-wires.
\end{enumerate}
Without property (2), a Prover could evaluate the DSGC with an
arbitrary $x_D'$ (using publicly-derivable input labels) and recover
$\msg$ using a proof for $(x_S, x_D')$, defeating the binding. Both
properties together are necessary for DSGC binding
(Lemma~\ref{lem:dsgc-binding}).

\subsection{Construction}

Let
\begin{equation}
  P_S \;=\; \sum_{j \in S} x_j \cdot L_j,\qquad
  P_D \;=\; \sum_{j \in D} x_j \cdot L_j,\qquad
  Y_S \;=\; e(\alpha, \beta) \cdot e(P_S, \gamma).
  \label{eq:ys-def}
\end{equation}
By linearity of $vk_x$ and bilinearity of $e$:
\begin{equation}
  vk_x \;=\; P_S + P_D,\qquad
  Y \;=\; Y_S \cdot e(P_D, \gamma).
  \label{eq:y-split}
\end{equation}

\begin{algorithm}[H]
\caption{$\Enc^\star(vk, x_S, S, D;\, I)$ \hfill{\footnotesize Verifier, off-chain, $T_0$}}
\label{alg:pb-enc}
\begin{algorithmic}[1]
\Require $vk$; static values $x_S$; partition $(S, D)$;
  instance $I = (r,\, \msg,\, \Delta,\, dsgc\_ct,\, adaptor)$
  where $dsgc\_ct$ garbles Alg.~\ref{alg:dsgc}
\Ensure $ct = (dsgc\_ct,\, adaptor,\, ct_2,\, ct_3)$
\State $ct_2 \gets r \cdot \delta$
\State $P_S \gets \sum_{j \in S} x_j \cdot L_j$
\State $Y_S \gets e(\alpha, \beta) \cdot e(P_S, \gamma)$
\State $ct_3 \gets \msg \oplus H(r \cdot Y_S)$ \Comment{mask uses $Y_S$, not $Y$}
\State \Return $(dsgc\_ct,\, adaptor,\, ct_2,\, ct_3)$
\end{algorithmic}
\end{algorithm}

\begin{algorithm}[H]
\caption{$\mathsf{DSGC}$ functionality \hfill{\footnotesize garbled at $T_0$, evaluated at $T_2$}}
\label{alg:dsgc}
\begin{algorithmic}[1]
\Require $\pi_1 \in \Gone$ (Lamport-revealed); $x_D \in \Fr^{|D|}$ (Lamport-revealed)
  \Statex constants (baked at garbling): $r \in \Fr$, $\{r \cdot L_j\}_{j \in D} \subset \Gone$
\Ensure $(c_1, c_1') = (r \cdot \pi_1,\ r \cdot P_D) \in \Gone^2$
\State $c_1 \gets r \cdot \pi_1$ \Comment{variable-base scalar-mul}
\State $c_1' \gets \sum_{j \in D} x_j \cdot (r \cdot L_j)$ \Comment{$|D|$ fixed-base scalar-muls, summed}
\State \Return $(c_1, c_1')$
\end{algorithmic}
\end{algorithm}

\begin{algorithm}[H]
\caption{$\Dec^\star(\pi, x = (x_S, x_D), ct)$ \hfill{\footnotesize Prover, off-chain, $T_2$}}
\label{alg:pb-dec}
\begin{algorithmic}[1]
\Require $\pi = (\pi_1, \pi_2, \pi_3)$ for $(vk, (x_S, x_D))$; $ct$
\Ensure $\msg$
\State $(ct_1, ct_1') \gets \mathsf{DSGC}(\pi_1,\, x_D)$ \Comment{via GC unlock; both outputs together}
\State $Q \gets e(ct_1',\, \gamma)$ \Comment{$Q = e(P_D, \gamma)^r$}
\State $r \cdot Y \gets e(ct_1,\, \pi_2) / e(\pi_3,\, ct_2)$ \Comment{by \eqref{eq:g16}}
\State $r \cdot Y_S \gets (r \cdot Y) / Q$ \Comment{by \eqref{eq:y-split}}
\State $\msg \gets ct_3 \oplus H(r \cdot Y_S)$
\State \Return $\msg$
\end{algorithmic}
\end{algorithm}

\begin{lemma}[Correctness]\label{lem:correct}
Algorithms~\ref{alg:pb-enc} and \ref{alg:pb-dec} satisfy the
correctness condition of Definition~\ref{def:pbwe}, provided the
DSGC is garbled consistently with Algorithm~\ref{alg:dsgc} and the
Prover supplies Lamport-revealed $(\pi_1, x_D)$ matching the proof's
public input.
\end{lemma}

\begin{proof}
Let $\pi$ be valid for $(vk, (x_S, x_D))$. The DSGC on
$(\pi_1, x_D)$ outputs $ct_1 = r \cdot \pi_1$ and
$ct_1' = r \cdot P_D$. By \eqref{eq:g16},
$e(ct_1, \pi_2) / e(\pi_3, ct_2) = r \cdot Y$. Line~2 yields
$Q = e(P_D, \gamma)^r$. Applying \eqref{eq:y-split}:
$r \cdot Y / Q = r \cdot Y_S$. Since
$ct_3 = \msg \oplus H(r \cdot Y_S)$, line~5 recovers $\msg$.
\end{proof}

\subsection{On-chain surface}

Algorithms~\ref{alg:pb-enc}--\ref{alg:pb-dec} run entirely off-chain.
The on-chain surface of partial-binding equals that of basic BABE
plus Lamport commitments on $x_D$ in the Assert transaction (routed
into DSGC input wires, analogous to Lamport-on-$\pi_1$). No new
Bitcoin-script primitive is required; only additional Lamport slots
per bit of $x_D$.

% ============================================================
\section{Protocol integration: deferred-proof setting}\label{sec:protocol}

The partial-binding scheme is designed for the deferred-inputs
pattern of \S\ref{sec:prelim} where the Groth16 proof is generated
\emph{after} encryption. We describe the integration with the BitVM2
bridge dispute flow, which is the canonical instantiation.

\subsection{Timeline}

Let $T_0 < T_1 < T_2 < T_3$ denote the protocol's four observable
events on Bitcoin.

\begin{center}\small
\renewcommand{\arraystretch}{1.25}
\begin{tabular}{@{}l p{4.0cm} p{7.2cm}@{}}
\toprule
\textbf{Event} & \textbf{Role} & \textbf{Actions} \\
\midrule
$T_0$ (peg-in) & Verifier, off-chain & Compute $x_S$; run
$\Enc^\star(vk, x_S, S, D;\ I)$ for each C\&C instance; publish
commitments to $(ct_2, ct_3, ct_D)$ on-chain ($\blacksquare_{\mathrm{com}}$). \\
$T_1$ (Kickoff) & Prover, on-chain & Broadcast Kickoff tx; opens
challenge window. No proof yet. $x_D$ unknown. \\
$[T_1, T_2]$ & Prover, off-chain & Observe watchtower challenges;
select $x_D$ (e.g.\ current Bitcoin tip, watchtower bitmap); build
auxiliary proofs (header chain, state chain); run Groth16 prover to
produce $\pi = \mathsf{Prove}(vk, (x_S, x_D), w)$. \\
$T_2$ (Assert) & Prover, on-chain & Broadcast Assert tx with Lamport
reveals of $\pi_1$ and $x_D$; GC gate-by-gate unlocking yields
$ct_1 = r \cdot \pi_1$. Run $\Dec^\star(\pi, (x_S, x_D), ct)$
off-chain to obtain $\msg$. \\
$T_3$ (Take) & Prover, on-chain & Spend dispute UTXO via hashlock
$\mathrm{RIPEMD160}(\mathrm{SHA256}(\msg))$. Take1/Take2 scripts
enforce byte-equality between Lamport-revealed $x_D$ and the proof's
$x_D$. \\
\bottomrule
\end{tabular}
\end{center}

\subsection{Why partial-binding is required here}

At $T_0$, the Verifier does not know the Groth16 proof or the
runtime values $x_D$. Basic BABE cannot encrypt: its $\Enc$
(Alg.~\ref{alg:babe-enc}, line 2) requires $\sum_{j=0}^{n} x_j L_j$,
which needs every $x_j$. Postponing encryption until $T_2$ is not
viable because the Verifier's role ends at peg-in; keeping Verifier
online and interactive until dispute time contradicts the trust
model (Verifier committee may rotate, become unreachable, or be
assumed cryptographically absent after peg-in).

Partial-binding lifts exactly this constraint: $\Enc^\star$
(Alg.~\ref{alg:pb-enc}) needs only $x_S$, which is peg-in-knowable.
The runtime values are absorbed at decryption
(Alg.~\ref{alg:pb-dec}, lines 2--5) using the published $ct_D$.

\subsection{Consistency across the timeline}

Correctness and soundness across the full timeline reduce to four
composable properties, each discharged by a distinct mechanism.

\begin{center}\small
\renewcommand{\arraystretch}{1.25}
\begin{tabular}{@{}p{3.8cm} p{8cm}@{}}
\toprule
\textbf{Property} & \textbf{Enforced by} \\
\midrule
$x_S$ is peg-in-fixed & $ct_3$ binding via $Y_S$ in $\Enc^\star$
(Lemma~\ref{lem:pbwe}, prefix soundness). \\
$\pi$ is valid for $(x_S, x_D)$ & Lemma~\ref{lem:dsgc-binding}: the
DSGC binds $r \cdot P_D$ to Lamport-committed $x_D$; a mismatched
proof fails bilinear cancellation. \\
$x_D$ is unique (no equivocation) & Lamport commitment on $x_D$ in
Assert tx: double-reveal discloses secret and burns deposit. \\
$x_D$ used in $\pi$ matches Lamport-committed $x_D$ &
DSGC binding (\S\ref{sec:dsgc}, Lemma~\ref{lem:dsgc-binding}):
$r \cdot P_D$ output is uniquely determined by Lamport-revealed
$x_D$; substitution attempts fail the bilinear cancellation. \\
\bottomrule
\end{tabular}
\end{center}

The four properties compose multiplicatively in the security
reduction (Lemma~\ref{lem:pbwe}, Lemma~\ref{lem:dsgc-binding}).

\subsection{Failure modes and their resolutions}

\begin{center}\small
\renewcommand{\arraystretch}{1.25}
\begin{tabular}{@{}p{5.5cm} p{6.5cm}@{}}
\toprule
\textbf{Failure} & \textbf{Resolution} \\
\midrule
Prover produces $\pi$ with wrong $x_S$ & Verifier equation
cancellation in $\Dec^\star$ line 5 fails; $\msg$ not recovered;
dispute path blocked. \\
Prover chooses $x_D$ not consistent with on-chain observations
(e.g.\ fails watchtower dual-check) & Circuit-internal constraints
reject witness; $\mathsf{Prove}$ has no output; Prover cannot
construct $\pi$. \\
Prover attempts to reveal two distinct $x_D$ values & Lamport
equivocation exposes secret key; Verifier slashes deposit. \\
Prover generates $\pi$ too late for the challenge window & Protocol
liveness failure, not security. Operator liveness parameters must
size the window against Groth16 proving time. \\
Verifier garbles DSGC incorrectly & BABE C\&C soundness: the opened
instance subset is re-garbled from seed and checked for equality;
Prover aborts before deposit if any opened instance disagrees. \\
\bottomrule
\end{tabular}
\end{center}

\subsection{Protocol invariant}

\begin{theorem}[Deferred-proof correctness]\label{thm:deferred}
Under Lemmas~\ref{lem:correct}, \ref{lem:pbwe},
\ref{lem:dsgc-binding}, together with Lamport one-wayness and
cut-and-choose soundness inherited from BABE, the following holds:
for every ciphertext $ct \gets \Enc^\star(vk, x_S, S, D;\ I)$
produced at $T_0$ and every PPT Prover strategy over $[T_0, T_3]$,
the Prover obtains $\msg$ if and only if they commit on-chain to a
specific $x_D^\dagger$ at $T_2$ and produce a valid Groth16 proof
$\pi^\dagger$ for $(x_S, x_D^\dagger)$.
\end{theorem}

\begin{proof}[Proof sketch]
$(\Leftarrow)$ is Lemma~\ref{lem:correct}.
$(\Rightarrow)$ follows by Lemma~\ref{lem:dsgc-binding}: absent a
$\pi^\dagger$ consistent with Lamport-committed $x_D^\dagger$, the
bilinear cancellation in $\Dec^\star$ fails. Absent the committed
$x_S$, the $Y_S$ factor mismatches (Lemma~\ref{lem:pbwe}).
\end{proof}

% ============================================================
\section{Security analysis}\label{sec:security}

\subsection{Main theorem}

\begin{lemma}[Partial-binding security]\label{lem:pbwe}
Let $\mathcal{A}$ be a PPT adversary against the partial-binding WE
scheme defined by Algorithms~\ref{alg:pb-enc}--\ref{alg:pb-dec}.
There exists a PPT reduction such that
\[
  \Adv^{\mathrm{pbWE}}_{\mathcal{A}}(\lambda)
  \;\leq\;
  \Adv^{\mathrm{BABE}}(\lambda) \,+\, \Adv^{\mathrm{DLog}}_{\Gone}(\lambda).
\]
\end{lemma}

\begin{proof}[Proof sketch]
The scheme differs from basic BABE in two places: the single-output
scalar-mul GC is replaced with the DSGC of
Algorithm~\ref{alg:dsgc}, and the mask in $ct_3$ uses $Y_S$ rather
than $Y$. GC security bounds the probability that an adversary
learns unintended outputs of the DSGC evaluation. The $Y$-to-$Y_S$
shift is reconciled via the bilinearity identity \eqref{eq:y-split}
inside the decryption equation. Any attack that recovers $\msg$
without a valid proof extracts proof elements from the GC
transcript, yielding an attack on basic BABE.
\end{proof}

\subsection{$r$-leakage analysis}

\begin{lemma}[No new $r$-leakage]\label{lem:r-leak}
Let $\mathcal{B}$ be a PPT adversary given the artifacts published
by $\Enc^\star$ plus the dispute-time DSGC outputs $(ct_1, ct_1')$.
Then
\[
  \Adv^{r\text{-leak}}_{\mathcal{B}}(\lambda)
  \;\leq\;
  \Adv^{\mathrm{DLog}}_{\Gone}(\lambda) \,+\, \Adv^{\mathrm{coDLog}}(\lambda).
\]
\end{lemma}

\begin{proof}[Proof sketch]
Public exposures of $r$-scaled samples are $ct_2 = r \cdot \delta$,
$ct_1 = r \cdot \pi_1$, and $ct_1' = r \cdot P_D$. These are
multi-instance DLog challenges with a shared scalar; self-reducibility
reduces to single-instance DLog. Cross-group combinations add at
most a co-DLog term.
\end{proof}

\subsection{$x_D$-substitution attack on a naive scheme}\label{sec:subst}

A na\"{i}ve instantiation of partial-binding WE that publishes
$ct_D = \{r \cdot L_j\}_{j \in D}$ and lets the decryptor form
$r \cdot P_D$ off-chain from $x_D$ and $ct_D$ admits the following
attack:

\begin{lemma}[Substitution attack on naive $ct_D$ scheme]\label{lem:attack}
Let $\mathcal{A}$ control a Prover possessing a valid proof $\pi'$
for $(x_S, x_D')$, and let $x_D^\dagger \neq x_D'$ be the value that
enclosing-protocol commitments (Lamport reveal, chain-derived
bitmap) require. Against the naive $ct_D$ instantiation, $\mathcal{A}$
recovers $\msg$ via $\Dec^\star$ by (i) Lamport-committing
$x_D^\dagger$ on-chain and (ii) using $x_D'$ off-chain to form
$r \cdot P_{D'}$. No on-chain check discriminates between
$x_D^\dagger$ (committed) and $x_D'$ (used in decryption) because
the proof's public-input vector is not script-verifiable.
\end{lemma}

\begin{proof}[Proof sketch]
In the naive scheme, $r \cdot P_D$ is Prover-computed from public
$ct_D$ and arbitrary $x_D$; the verifier equation cancellation
succeeds for $(x_S, x_D', \pi')$. Lamport on $x_D^\dagger$ pins
\emph{that} value for on-chain equality checks but provides no link
to the off-chain computation of $r \cdot P_D$.
\end{proof}

\subsection{DSGC binding closes the attack}

The construction of \S\ref{sec:construction} replaces $ct_D$ with the
DSGC of Algorithm~\ref{alg:dsgc}: $r \cdot P_D$ becomes a GC output
parameterized by Lamport-committed $x_D$, rather than a Prover-chosen
off-chain quantity.

\begin{lemma}[DSGC binding]\label{lem:dsgc-binding}
Let $\mathcal{A}$ control a Prover who produces $\msg$ via
$\Dec^\star$ instantiated with Algorithm~\ref{alg:dsgc}. Let
$x_D^{\mathrm{Lam}}$ denote the Lamport-revealed dynamic value and
$x_D^{\pi}$ the public-input vector of $\mathcal{A}$'s proof.
Then $x_D^{\mathrm{Lam}} = x_D^{\pi}$ except with probability
\[
  \Adv^{\mathrm{BABE}}(\lambda) \,+\, \Adv^{\mathrm{GC\text{-}sec}}(\lambda)
  \,+\, \Adv^{\mathrm{Lamport}}(\lambda) \,+\, \Adv^{\mathrm{G16\text{-}KS}}(\lambda).
\]
\end{lemma}

\begin{proof}[Proof sketch]
Two facts, each reducing to a named assumption:

\emph{(a) GC privacy hides $\{r \cdot L_j\}_{j \in D}$.} Constants
baked into the DSGC at garbling time are not extractable by the
evaluator: by GC privacy, the evaluator's view on a single
evaluation reveals only the output labels for the specific inputs
provided, not constants or intermediate wire values
(Lindell--Pinkas; Bellare--Hoang--Rogaway). Hence the adversary
cannot directly compute $r \cdot P_{D'}$ for $x_D' \neq x_D^{\mathrm{Lam}}$
by combining extracted $\{r \cdot L_j\}$.

\emph{(b) Lamport soldering fixes DSGC input to $x_D^{\mathrm{Lam}}$.}
By construction, the DSGC input-wire labels for $x_D$ are derived
from Lamport preimages ($L_{j,i}^b = \mathrm{KDF}(s_{j,i}^b)$). To
evaluate the DSGC with some $x_D'$, the Prover needs labels
$\{L_{j,i}^{b'}\}$ for $x_D'$'s bit pattern. Revealing preimages for
a bit pattern different from the on-chain Lamport reveal constitutes
equivocation: it simultaneously exposes both halves of at least one
Lamport slot, which is slashable.

Combining (a) and (b): any DSGC evaluation accessible to the Prover
yields $r \cdot P_D$ with $x_D = x_D^{\mathrm{Lam}}$. Therefore line~2
of Alg.~\ref{alg:pb-dec} computes
$Q = e(P_{D, \mathrm{Lam}}, \gamma)^r$. The decryption quotient
$r \cdot Y / Q$ evaluates to $r \cdot Y_S$ iff
$x_D^{\pi} = x_D^{\mathrm{Lam}}$; otherwise the residual factor
$e(P_{D, \pi} - P_{D, \mathrm{Lam}}, \gamma)^r \neq 1$ and the mask
$H(r \cdot Y_S)$ is not recovered. A Prover who still outputs $\msg$
either inverts $H$ or produces a proof extracting to no witness,
contradicting Groth16 KS.
\end{proof}

\begin{remark}[Proactive vs.\ reactive binding]
The DSGC mechanism binds $x_D$ proactively through the garbled
circuit's structure: substitution is infeasible without breaking GC
privacy or Lamport one-wayness, and no Challenger is required. In
contrast, a disprove-based design would rely on an active Challenger
running the full Groth16 verifier and posting a Disprove transaction
upon detecting $x_D^{\pi} \neq x_D^{\mathrm{Lam}}$, reintroducing
full-verifier cost on the dispute path. DSGC preserves BABE's core
``no full verifier on chain'' property.
\end{remark}

\begin{remark}[Both GC privacy and soldering are required]
Lemma~\ref{lem:dsgc-binding} depends on \emph{both} assumptions
jointly. Weakening either enables the attack: if the Verifier
published $\{r \cdot L_j\}$ in plaintext (violating GC privacy) the
Prover forms $r \cdot P_{D'}$ directly; if the DSGC input wires for
$x_D$ are not soldered to Lamport preimages (allowing public label
derivation) the Prover evaluates the DSGC on any $x_D'$ of their
choice. The construction therefore requires that $\{r \cdot L_j\}$
be baked as garbled constants \emph{and} that $x_D$ input wires be
properly soldered. Both are standard ingredients of BABE's GC
infrastructure applied to the extended input space.
\end{remark}

% ============================================================
\section{Cost}\label{sec:cost}

\begin{center}
\small
\renewcommand{\arraystretch}{1.2}
\begin{tabular}{@{}l l l@{}}
\toprule
\textbf{Item} & \textbf{Basic BABE} & \textbf{Partial-binding (delta)} \\
\midrule
GC: output count                        & $r \cdot \pi_1$             & $+\,r \cdot P_D$ (DSGC) \\
GC: gate count                          & $\sim 4.7 \times 10^8$      & $+\,\sim |D| \cdot 1.5 \times 10^8$ (fixed-base) \\
Encryption: $\Gone$ scalar-muls         & $O(n)$                      & $+\,|D|$ (baked into DSGC constants) \\
Encryption: pairings                    & $2$                         & unchanged \\
Ciphertext size                         & baseline                    & unchanged \\
Decryption: $\Gone$ scalar-muls         & --                          & $+\,|D|$ (inside DSGC evaluation) \\
Decryption: pairings                    & $2$                         & $+\,1$ \\
Decryption: $\Gt$ divisions             & $1$                         & $+\,1$ \\
Lamport slots on-chain                  & $\pi_1$ only (\textasciitilde 508 bits) & $+\,|D| \cdot 254$ bits for $x_D$ \\
On-chain primitives                     & hashlock, GC, Lamport       & unchanged \\
\bottomrule
\end{tabular}
\end{center}

For the BitVM2 bridge with $|D| = 4$ field slots: DSGC grows to
$\sim\!1.1 \times 10^9$ gates (roughly $2.3\times$ baseline),
dispute-path Taproot-unlock witness grows proportionally, and the
Lamport witness at Assert grows by $\sim\!10^3$ bits. The happy path
remains cheap: the Prover evaluates DSGC off-chain to recover $\msg$
and spends via hashlock. Full on-chain GC unlocking runs only in the
rare dispute path.

% ============================================================
\section{Conclusion}

Partial-binding WE resolves the full-public-input requirement of
basic BABE for protocols with deferred public inputs. The
construction uses a Dual-Scalar Garbled Circuit (\S\ref{sec:dsgc})
whose two outputs---$r \cdot \pi_1$ and $r \cdot P_D$---together
provide prefix binding on $x_S$ and proactive cryptographic binding
on $x_D$. The reduction stays within standard BABE assumptions plus
GC-security and Lamport one-wayness. No new on-chain primitive is
introduced; BABE's core property of avoiding a full on-chain Groth16
verifier is preserved.

% ============================================================
\appendix
\section{Universal-Composability analysis}\label{app:uc}

We sketch a UC treatment of the partial-binding WE scheme. A full
simulator-based proof is deferred; this appendix fixes the
functionality, the hybrid model, and the simulation outline
sufficient for a reviewer to verify the security architecture.

\subsection{Ideal functionality $\mathcal{F}_{\mathrm{pbWE}}$}

$\mathcal{F}_{\mathrm{pbWE}}$ interacts with Verifier $V$,
Prover $P$, and adversary $\mathcal{S}$. Internal state is a
per-session record indexed by $\mathsf{sid}$.

\begin{center}
\begin{tabular}{@{}p{15cm}@{}}
\toprule
\textbf{Setup (at $T_0$):} Upon $(\mathsf{ENC}, \mathsf{sid}, vk, x_S, S, D, \msg)$ from $V$: \\
\quad store $(\mathsf{sid}, vk, x_S, S, D, \msg, \bot)$; send
$(\mathsf{ENCRYPTED}, \mathsf{sid}, x_S, S, D)$ to $\mathcal{S}$ and
all parties. \emph{$\msg$ is not leaked to $\mathcal{S}$.} \\
\midrule
\textbf{Commit (at $T_2$):} Upon $(\mathsf{COMMIT}, \mathsf{sid}, x_D, h_{\pi_1})$ from $P$: \\
\quad if previously committed, reject (equivocation); else store
$\mathsf{committed} := (x_D, h_{\pi_1})$ and send
$(\mathsf{COMMITTED}, \mathsf{sid}, x_D, h_{\pi_1})$ to $\mathcal{S}$. \\
\midrule
\textbf{Decrypt (at $T_2$):} Upon $(\mathsf{DEC}, \mathsf{sid}, \pi)$ from $P$: \\
\quad require prior commit $(x_D^*, h^*)$; return $(\mathsf{DEC\_OK}, \mathsf{sid}, \msg)$ \\
\quad iff $\mathsf{Verify}_{\mathrm{G16}}(vk, (x_S, x_D^*), \pi) = 1$
and $H(\pi.\pi_1) = h^*$; else $(\mathsf{DEC\_FAIL}, \mathsf{sid})$. \\
\bottomrule
\end{tabular}
\end{center}

The functionality encodes four security properties by construction:
prefix binding (on $x_S$), proactive $x_D$ binding (at COMMIT),
decryption conditional on valid proof matching the committed
$(x_S, x_D)$, and equivocation prevention.

\subsection{Hybrid model and trust setup}

The real protocol is analyzed in the
$(\mathcal{F}_{\mathrm{CRS}}, \mathcal{F}_{\mathrm{BTC}}, \mathcal{F}_{\mathrm{RO}})$-hybrid model:
\begin{itemize}
  \item $\mathcal{F}_{\mathrm{CRS}}$: ideal Groth16 trusted setup delivering $vk$.
  \item $\mathcal{F}_{\mathrm{BTC}}$: idealized Bitcoin ledger exposing transaction inclusion, Taproot unlock, and Lamport preimage reveal.
  \item $\mathcal{F}_{\mathrm{RO}}$: random oracle modeling $H$.
\end{itemize}

\subsection{Simulator outline}

The simulator $\mathcal{S}$ plays the real-world adversary's role
against $\mathcal{F}_{\mathrm{pbWE}}$:

\begin{itemize}
  \item \emph{At ENC.} $\mathcal{S}$ learns $(x_S, S, D)$ but not
    $\msg$. Samples $\widetilde{\msg}$ uniformly and runs
    $\Enc^\star$ honestly with $\widetilde{\msg}$, producing a
    simulated ciphertext $\widetilde{ct}$. Indistinguishability
    reduces to partial-binding security
    (Lemma~\ref{lem:pbwe}).
  \item \emph{At COMMIT.} $\mathcal{S}$ extracts
    $(x_D, \pi_1)$ from the Prover's Lamport reveals (Lamport
    one-wayness gives a unique extraction). Forwards
    $(x_D, H(\pi_1))$ to $\mathcal{F}_{\mathrm{pbWE}}$.
  \item \emph{At DEC.} $\mathcal{S}$ applies the Groth16 KS extractor
    to $\pi$, forwards $\pi$ to $\mathcal{F}_{\mathrm{pbWE}}$. On
    $(\mathsf{DEC\_OK}, \msg)$, $\mathcal{S}$ programs
    $\mathcal{F}_{\mathrm{RO}}$ so that $\widetilde{ct}_3 \oplus H(r \cdot Y_S) = \msg$,
    preserving transcript consistency. On
    $(\mathsf{DEC\_FAIL})$, $\mathcal{S}$ emits the residual garbage
    that $\Dec^\star$ would compute (Lemma~\ref{lem:dsgc-binding}).
\end{itemize}

\subsection{Hybrid argument}

The real-ideal indistinguishability is established by a chain of
hybrids, each justified by a standard assumption:

\begin{center}\small
\renewcommand{\arraystretch}{1.15}
\begin{tabular}{@{}l l@{}}
\toprule
\textbf{Hybrid step} & \textbf{Justification} \\
\midrule
Real $\pi_{\mathrm{pbWE}}$ & --- \\
$ct_3$ masked by RO output & programmable $\mathcal{F}_{\mathrm{RO}}$ \\
DSGC transcript simulated with random labels & GC-security of Alg.~\ref{alg:dsgc} \\
Lamport reveals simulated via trapdoor & Lamport one-wayness \\
Groth16 $\pi$ extracted & Groth16 KS in AGM \\
$\widetilde{\msg}$ replaced with real $\msg$ & $\mathcal{F}_{\mathrm{pbWE}}$ on $\mathsf{DEC\_OK}$ \\
Ideal $\mathcal{F}_{\mathrm{pbWE}}$ & --- \\
\bottomrule
\end{tabular}
\end{center}

Each step's distinguishing advantage is bounded by a negligible
term in $\lambda$.

\subsection{UC statement (informal)}

\begin{theorem}[UC emulation]\label{thm:uc}
Under the cryptographic assumptions of \S\ref{sec:security} and the
hybrid model above, there exists a PPT simulator $\mathcal{S}$ such
that for every PPT environment $\mathcal{Z}$ and PPT real-world
adversary $\mathcal{A}$,
\[
  \bigl|\Pr[\mathsf{EXEC}^{\pi_{\mathrm{pbWE}}}_{\mathcal{A}, \mathcal{Z}} = 1]
       - \Pr[\mathsf{EXEC}^{\mathcal{F}_{\mathrm{pbWE}}}_{\mathcal{S}, \mathcal{Z}} = 1]\bigr|
  \;\leq\; \negl(\lambda).
\]
\end{theorem}

\subsection{Composition consequences}

UC emulation yields:
\begin{itemize}
  \item \emph{Sequential composition.} Multiple peg-in sessions
    (distinct $\mathsf{sid}$) remain secure simultaneously.
  \item \emph{Subroutine replacement.} Replacing Groth16 by any
    UC-secure SNARK preserves $\mathcal{F}_{\mathrm{pbWE}}$
    emulation, modulo the replacement's own UC proof.
  \item \emph{Enclosing-protocol compatibility.} Partial-binding WE
    composes with adaptor signatures, Taproot disprove flows, and
    higher-layer bridge protocols without rerunning its security
    proof.
\end{itemize}

A full mechanized UC proof is future work; \S\ref{sec:security} gives
the game-based analogue used for concrete parameters.

% ============================================================
\begin{thebibliography}{9}
\bibitem{babe2026}
Goat Research Team.
\newblock BABE: Verifying Proofs on Bitcoin Made 1000$\times$ Cheaper.
\newblock IACR ePrint 2026/065.

\bibitem{groth16}
J.~Groth.
\newblock On the Size of Pairing-Based Non-interactive Arguments.
\newblock EUROCRYPT 2016, pp.\ 305--326.

\bibitem{canetti-uc}
R.~Canetti.
\newblock Universally Composable Security: A New Paradigm for
Cryptographic Protocols.
\newblock FOCS 2001, pp.\ 136--145; updated in IACR ePrint 2000/067.
\end{thebibliography}

\end{document}
